% ==============================================================================
% INTRODUCCIÓN
% ==============================================================================
% Instrucciones de la plantilla:
% - Proporcionar un contexto claro sobre los objetivos, fundamentos teóricos y 
%   relevancia de la práctica realizada.
% - Incluir explicación de los conceptos gráficos abordados en la práctica 
%   (ej: rasterización, transformaciones geométricas, algoritmos de dibujo, 
%   OpenGL/WebGL, shaders, etc.). Se recomienda usar la información investigada 
%   en los previos para elaborar esta sección.
% - ES OBLIGATORIO incluir citas a fuentes bibliográficas en formato APA.
%   Ejemplo: "De acuerdo con \cite{teodoro_vite}, el cómputo gráfico..."
% ==============================================================================

\section{Introducción}

El cómputo gráfico moderno se fundamenta...

% Continuar la redacción integrando:
% 1. Contexto y motivación del tema.
% 2. Conceptos teóricos clave aplicados (rasterización, shaders, matrices MVP, etc.).
% 3. Citas bibliográficas correspondientes en formato APA.
